
\documentclass[11pt,a4paper]{article}
\usepackage[a4paper,margin=2.2cm]{geometry}
\usepackage[T1]{fontenc}
\usepackage[utf8]{inputenc}
\usepackage[spanish,es-noquoting]{babel}
\usepackage{lmodern}
\usepackage{microtype}
\usepackage{amsmath,amssymb,mathtools}
\usepackage{booktabs,longtable,array}
\usepackage{xcolor}
\usepackage{hyperref}
\hypersetup{colorlinks=true,linkcolor=black,urlcolor=blue}
\setlength{\parindent}{0pt}
\setlength{\parskip}{0.72em}
\title{\bfseries HMT -- UNIDAD--EULER--01\\Unidad, acumulaci\'on, pliegue y sombra tri\'adica de Euler}
\author{Ventana de investigaci\'on HMT}
\date{}
\begin{document}
\maketitle

\section*{Prop\'osito}
Esta ventana recoge una l\'inea conceptual que no debe perderse en la acumulaci\'on de c\'alculos: la unidad no aparece al final como un simple n\'umero neutro, sino como una frontera que se parte, se recompone y se transmite. La tesis operativa es:
\[
\boxed{\text{la suma acumula; el producto pliega; la unidad aparece como cierre de ambas lecturas.}}
\]

En esta ventana se auditan cinco hechos elementales pero estructuralmente fuertes: la descomposici\'on
\(1000=729+271\), la identidad \(196884=54(3645+1)\), la sombra exacta
\((90-120)/360=-1/12\), la lectura espectral del canal \(90/120\) en la rueda SU--12, y la condici\'on de unidad de \(\varphi\), donde sumar una unidad y multiplicar por \(\varphi\) son la misma operaci\'on estructural.

\section{La unidad como bulk m\'as corona}
Una triada decimal vive en el cubo completo
\[
\{0,\ldots,9\}^3,
\]
con cardinal
\[
10^3=1000.
\]
El bulk activo, sin ceros, es
\[
\{1,\ldots,9\}^3,
\qquad |\{1,\ldots,9\}^3|=9^3=729.
\]
La corona decimal es el complemento:
\[
1000-729=271.
\]
Por tanto:
\[
\boxed{1000=729+271.}
\]
Normalizando por la unidad decimal:
\[
\boxed{1=0.729+0.271.}
\]
Esta igualdad no debe leerse como una curiosidad decimal. En HMT significa que la unidad observable en base--1000 se descompone en bulk tri\'adico activo y corona de lectura. La unidad no es homog\'enea: tiene interior y borde.

\section{Suma como acumulaci\'on y producto como pliegue}
En APP bidimensional definimos dos lectores:
\[
APP_+(i,j)=dr_9(i+j),
\qquad
APP_\times(i,j)=dr_9(ij).
\]
La suma desplaza y acumula. El producto pliega e identifica. La auditor\'ia finita da:
\[
S_{+,2}=405,
\qquad
S_{\times,2}=459,
\]
por lo que
\[
\boxed{\Delta_{2D}=S_{\times,2}-S_{+,2}=54.}
\]
El mismo plano produce
\[
\Sigma(9)=S_{+,2}+S_{\times,2}=864,
\]
\[
1728=2\Sigma(9),
\qquad
744=\Sigma(9)-120.
\]
As\'i, APP no es una tabla: es una superficie de doble lectura. La diferencia entre acumulaci\'on y pliegue produce un defecto entero.

\section{APP c\'ubica y el ancla 196884}
En la APP c\'ubica activa se obtiene:
\[
S_{+,3}=3645,
\qquad
S_{\times,3}=4617.
\]
El dato estructural central es
\[
\boxed{S_{+,3}=3645.}
\]
Combinando el defecto bidimensional con el bulk aditivo tridimensional:
\[
\boxed{196884=54(3645+1).}
\]
Esta identidad es la forma aritm\'etica HMT de una puerta Moonshine. No se invoca aqu\'i el Monstruo como hip\'otesis; aparece primero la aritm\'etica APP:
\[
\text{defecto APP 2D}\times(\text{bulk APP 3D}+1).
\]
El \(+1\) es esencial: en el lenguaje Moonshine cl\'asico, el primer coeficiente grande se separa como \(196883+1\). En HMT, ese uno debe interpretarse como unidad trivial, cierre o modo constante que acompa\~na a la simetr\'ia no trivial.

\section{Euler tri\'adico: la sombra exacta de \(-1/12\)}
La relaci\'on \(90/120\) tiene una sombra normalizada exacta:
\[
\boxed{\frac{90-120}{360}=-\frac{1}{12}.}
\]
Y en la lectura hexada--octada:
\[
90=6\binom{6}{2},
\qquad
120=8\binom{6}{2},
\qquad
360=24\binom{6}{2},
\]
por tanto:
\[
\boxed{\frac{6-8}{24}=-\frac{1}{12}.}
\]
Este es el punto que debe llamarse, con cuidado, \emph{sombra tri\'adica de Euler}. No se afirma que esta identidad por s\'i sola derive toda la regularizaci\'on de Euler/Riemann. Lo que se afirma es m\'as concreto y m\'as fuerte dentro de HMT: el defecto normalizado entre hexada visible y octada ambiente produce exactamente el mismo n\'umero finito que aparece en la regularizaci\'on \(\zeta(-1)=-1/12\).

Este dato conecta tres capas:
\[
\boxed{90/120}\quad\leftrightarrow\quad
\boxed{6/8}\quad\leftrightarrow\quad
\boxed{-1/12}.
\]
Por eso el canal \(90/120\) no debe tratarse como un adorno angular. Es la sombra entre centro y borde, entre hexada y octada, entre suma visible y producto plegado.

\section{Lectura espectral SU--12 del canal 90/120}
La rueda SU--12 permite representar el canal 90/120 por el patr\'on
\[
p_m=12\cos\left(\frac{2\pi\cdot 3m}{12}\right)-\cos\left(\frac{2\pi\cdot4m}{12}\right).
\]
Su transformada discreta de Fourier se concentra en los modos
\[
k=3,9
\qquad\text{y}\qquad
k=4,8.
\]
La raz\'on de amplitudes es
\[
\boxed{\frac{|\widehat p_3|}{|\widehat p_4|}=12.}
\]
As\'i, el canal 90/120 tiene una lectura espectral: excita el par de modos 3/4 con raz\'on 12:1. Esta observaci\'on conecta la rueda de doce fases con la sombra \(-1/12\): no se trata de una relaci\'on meramente decimal, sino de una relaci\'on de modos.

\section{La unidad \(\varphi\): sumar y multiplicar}
La raz\'on \'aurea satisface
\[
\varphi^2=\varphi+1,
\]
por lo que
\[
\boxed{\varphi+1=\varphi^2.}
\]
Tambi\'en
\[
\varphi-1=\frac1\varphi.
\]
Esto significa que, en la regi\'on \(\varphi\), sumar una unidad equivale a multiplicar por \(\varphi\), y restar una unidad equivale a dividir por \(\varphi\). Dicho en lenguaje HMT:
\[
\boxed{\text{en }\varphi,\text{ acumulaci\'on y escala se identifican por la unidad.}}
\]
Esta es la raz\'on por la que \(\varphi\) es mucho m\'as que una constante de proporci\'on. Es la regi\'on donde cantidad, escala, parte y todo se vuelven equivalentes bajo el cierre de la unidad.

\section{Euler producto: suma y producto como dos gram\'aticas de la unidad}
La identidad cl\'asica de Euler para la zeta de Riemann, en el dominio convergente, es
\[
\sum_{n\ge1}\frac{1}{n^s}
=
\prod_{p\ \mathrm{primo}}\left(1-p^{-s}\right)^{-1}.
\]
Esta ecuaci\'on es hist\'oricamente una de las grandes traducciones entre suma y producto. En HMT su papel conceptual es claro: la suma sobre todos los enteros y el producto sobre los primos son dos modos de leer la misma unidad aritm\'etica. La suma acumula t\'erminos; el producto pliega factores primarios.

Esta identidad cl\'asica no se usa aqu\'i como prueba nueva de HMT. Se usa como antecedente hist\'orico: la dualidad suma/producto ya estaba en el coraz\'on de la teor\'ia de n\'umeros. HMT propone una versi\'on finita de esa dualidad en APP:
\[
\boxed{\text{suma}=\text{acumulaci\'on},\qquad \text{producto}=\text{pliegue}.}
\]

\section{Lectura general}
La unidad aparece en varios niveles:
\begin{itemize}
\item Como \(1=0.729+0.271\): bulk activo m\'as corona decimal.
\item Como \(+1\) en \(196884=196883+1\): modo trivial unido a simetr\'ia no trivial.
\item Como cierre \(\varphi+1=\varphi^2\): unidad que convierte acumulaci\'on en escala.
\item Como sombra \(-1/12\): diferencia normalizada entre hexada y octada.
\item Como modo constante en la rueda SU--12.
\end{itemize}

La tesis compacta de esta ventana es:
\[
\boxed{\text{la unidad HMT no es un escalar pasivo; es el borde que permite pasar de acumulaci\'on a pliegue.}}
\]

\section*{Dictamen}
\[
\boxed{\text{Verde: }1000=729+271\text{ como partici\'on bulk/corona de la unidad.}}
\]
\[
\boxed{\text{Verde: }54\text{ como defecto suma/producto de APP 2D.}}
\]
\[
\boxed{\text{Verde: }196884=54(3645+1).}
\]
\[
\boxed{\text{Verde: }(90-120)/360=(6-8)/24=-1/12.}
\]
\[
\boxed{\text{Verde: lectura espectral 90/120 en modos 3/4 con raz\'on }12:1.}
\]
\[
\boxed{\text{Amarillo: integrar esta unidad tri\'adica con la emisi\'on profunda de fronteras }\Sigma_t.}
\]

\end{document}
